\section{考察と適用範囲}
写真の必須欄条件では，通常の別呼出し要約に対して成功と総token削減が両立した．しかし，実際に要約を使ったもう１件のCSVでは，入力が短くなっても成功を失った．要約の生成，入力への採用，総費用，最終成果を別々に測る必要がある．要約の作成費用だけを除いた比較や，成功した課題だけの合計は，全条件の性能維持を示さない．

対象は開発で使用した既知４課題，１モデル，各条件１回である．写真の7.50\%はこの実行対の観測値であり，未見課題の平均削減率や統計的な品質の非劣性ではない．３群は前の結果を見て設計した順次開発である．要約の作成時点，温度，指示，schema，出力配分と後続行動の差を含む方式全体の比較であり，同時生成だけの効果を分離していない．

全24枠の入力再構成も，実際の行動列と通知を固定した事後診断である．仮想入力を使った再生成と評価はしておらず，失われた意味の十分性を検証していない．同じseedでもcacheや数値実行，ホスト負荷の影響が残り，論理tokenを課金，FLOPs，実時間と同一視しない．完全履歴，課題データ，私有推論を公開せず，全条件の集約，固定条件，ソースと監査記録を公開する\footnote{\url{https://github.com/sophiria1225/AI-}}．
